\documentclass{article}
\usepackage{geometry}
\geometry{margin=1.5cm, vmargin={1.2cm,1.2cm}}

\usepackage{amsmath,amssymb,amsthm,amsfonts}
\usepackage{enumerate}
\usepackage{hyperref}

\newcommand{\C}{\mathbb{C}}
\newcommand{\R}{\mathbb{R}}
\newcommand{\Z}{\mathbb{Z}}
\newcommand{\dd}{\mathrm{d}}
\newcommand{\Res}{\operatorname{Res}}
\newcommand{\PV}{\operatorname{p.v.}}
\renewcommand{\Re}{\operatorname{Re}}
\renewcommand{\Im}{\operatorname{Im}}

\newenvironment{solution}{%
  \par\medskip\noindent\textbf{Solution.}\par
}{\par\medskip}

\title{\textbf{MIT Complex Analysis Recommended Exercises}}
\author{}
\date{}

\begin{document}
\maketitle

\noindent
These exercises are selected and retypeset from MIT OpenCourseWare
\href{https://ocw.mit.edu/courses/18-04-complex-variables-with-applications-spring-2018/}{18.04 Complex Variables with Applications, Spring 2018}.
The source problem sets are:
\href{https://ocw.mit.edu/courses/18-04-complex-variables-with-applications-spring-2018/mit18_04_s18_pset04.pdf}{PS4},
\href{https://ocw.mit.edu/courses/18-04-complex-variables-with-applications-spring-2018/mit18_04_s18_pset06.pdf}{PS6},
\href{https://ocw.mit.edu/courses/18-04-complex-variables-with-applications-spring-2018/mit18_04_s18_pset07.pdf}{PS7},
\href{https://ocw.mit.edu/courses/18-04-complex-variables-with-applications-spring-2018/mit18_04_s18_pset08.pdf}{PS8}, and
\href{https://ocw.mit.edu/courses/18-04-complex-variables-with-applications-spring-2018/mit18_04_s18_pset09.pdf}{PS9}.

\section{Priority Problems}

\subsection*{18.04 PS9, Problem 2}

Consider
\[
  f(z)=z^5-2z.
\]
\begin{enumerate}[(a)]
  \item How many times does $f$ wind the circle $|z|=3$ around the origin?
  That is, let $\gamma(\theta)$ parametrize the circle. How many times does
  the curve $f\circ\gamma(\theta)$ wind around the origin?

  \item How many times does $f$ wind the circle $|z|=1$ around the origin?

  \item How many times does $f$ wind the circle $|z|=3$ around the point
  $z=-2$?
\end{enumerate}

\begin{solution}
% Write your solution here.
\end{solution}

\subsection*{18.04 PS9, Problem 3}

\begin{enumerate}[(a)]
  \item Show that
  \[
    f(z)=z^3+9z+30
  \]
  has no roots in the disk $|z|<2$.

  \item Show that
  \[
    f(z)=z^6+4z^2-1
  \]
  has exactly two roots in the disk $|z|<1$.
\end{enumerate}

\begin{solution}
% Write your solution here.
\end{solution}

\subsection*{18.04 PS9, Problem 4}

Suppose $f(z)$ is analytic on a region containing $|z|\le 1$. Suppose also
that $|f(z)|<1$ on $|z|=1$. Show that $f(z)-z$ has exactly one zero in the
disk $|z|<1$.

\begin{solution}
% Write your solution here.
\end{solution}

\subsection*{18.04 PS6, Problem 4(b)}

Suppose $f$ is analytic and has a zero of order $m$ at $z_0$. Show that
\[
  g(z)=\frac{f'(z)}{f(z)}
\]
has a simple pole at $z_0$ with
\[
  \Res(g,z_0)=m.
\]

\begin{solution}
% Write your solution here.
\end{solution}

\subsection*{18.04 PS6, Problem 8}

Using variations of the geometric series, find the following series
expansions of
\[
  f(z)=\frac{1}{4-z^2}
\]
about $z_0=1$.
\begin{enumerate}[(a)]
  \item The Taylor series. What is the radius of convergence?

  \item The Laurent series on $1<|z-1|<R_1$. What is $R_1$?

  \item The Laurent series for $|z-1|>3$.
\end{enumerate}

\begin{solution}
% Write your solution here.
\end{solution}

\subsection*{18.04 PS6, Problem 5(b,c,d)}

\begin{enumerate}[(a)]
  \setcounter{enumi}{1}
  \item Find the residue of
  \[
    f_2(z)=\frac{z^2+1}{2z\cos z}
  \]
  at $z=0$.

  \item Let
  \[
    f_3(z)=\frac{e^z}{z(z+1)^3}.
  \]
  Find all the isolated singularities and compute the residue at each one.

  \item Find the residue at infinity of
  \[
    f_4(z)=\frac{1}{1-z}.
  \]
\end{enumerate}

\begin{solution}
% Write your solution here.
\end{solution}

\subsection*{18.04 PS6, Problem 6}

Write the principal part of each function at the isolated singularity.
Compute the corresponding residue.
\begin{enumerate}[(a)]
  \item
  \[
    f_1(z)=z^3e^{1/z}.
  \]

  \item
  \[
    f_2(z)=\frac{1-\cosh z}{z^3}.
  \]
\end{enumerate}

\begin{solution}
% Write your solution here.
\end{solution}

\subsection*{18.04 PS7, Problem 1(c)}

Compute
\[
  \int_0^{2\pi}\frac{\sin^2\theta}{a+b\cos\theta}\,\dd\theta,
  \qquad a>|b|>0.
\]
MIT gives the answer as
\[
  \frac{2\pi}{b^2}\left(a-\sqrt{a^2-b^2}\right).
\]

\begin{solution}
% Write your solution here.
\end{solution}

\subsection*{18.04 PS7, Problem 2(c)}

Show
\[
  \int_0^\infty \frac{1}{x^3+1}\,\dd x
  =
  \frac{2\pi\sqrt{3}}{9}
\]
by integrating around the boundary of a circular sector and letting
$R\to\infty$. The vertex angle of the sector is $2\pi/3$.

\begin{solution}
% Write your solution here.
\end{solution}

\subsection*{18.04 PS7, Problem 6}

Compute
\[
  \int_0^\infty\frac{\sqrt{x}}{x^2+1}\,\dd x.
\]
MIT gives the answer as
\[
  \frac{\pi}{\sqrt{2}}.
\]

\begin{solution}
% Write your solution here.
\end{solution}

\subsection*{18.04 PS8, Problem 2(b,c)}

\begin{enumerate}[(a)]
  \setcounter{enumi}{1}
  \item A fixed point $z$ of a transformation $T$ is one where $T(z)=z$.
  Let $T$ be a fractional linear transformation. Assume $T$ is not the
  identity map. Show $T$ has at most two fixed points.

  \item Let $S$ be a circle and $z_1$ a point not on the circle. Show that
  there is exactly one point $z_2$ such that $z_1$ and $z_2$ are symmetric
  with respect to $S$.

  Hint: start by proving this for $S$ a line.
\end{enumerate}

\begin{solution}
% Write your solution here.
\end{solution}

\subsection*{18.04 PS8, Problem 4}

\begin{enumerate}[(a)]
  \item Show that the mapping
  \[
    w=z+\frac{1}{z}
  \]
  maps the circle $|z|=a$, where $a\ne 1$, to the ellipse
  \[
    \frac{u^2}{(a+1/a)^2}+\frac{v^2}{(a-1/a)^2}=1.
  \]

  \item Where does it map the circle $|z|=1$?
\end{enumerate}

\begin{solution}
% Write your solution here.
\end{solution}

\subsection*{18.04 PS8, Problem 3}

Suppose you want to find a function $u$ harmonic on the right half-plane that
takes the values
\[
  u(0,y)=\frac{y}{1+y^2}
\]
on the imaginary axis. The first obvious guess is
\[
  u(z)=\Im\left(\frac{z}{1-z^2}\right).
\]
But this fails because $z/(1-z^2)$ has a singularity at $z=1$. Find a valid
$u$ using the following steps. Forget about this guess and go back to only
knowing that $u$ is harmonic and $u(0,y)=y/(1+y^2)$.

\begin{enumerate}[(a)]
  \item Show that rotation by $\alpha$ is a fractional linear transformation
  which corresponds to the matrix
  \[
    \begin{pmatrix}
      e^{i\alpha/2} & 0\\
      0 & e^{-i\alpha/2}
    \end{pmatrix}.
  \]

  \item Find a fractional linear transformation that maps the right half-plane
  to the unit disk, so that $u$ is transformed to a function $\phi$ with
  \[
    \phi(e^{i\theta})=\frac{\sin\theta}{2}.
  \]
  Hint: make sure $1$ is mapped to $0$. If your transformation still does not
  transform $u$ to the correct $\phi$, try composing with a rotation.

  \item Show that
  \[
    \phi(w)=\frac{1}{2}\Im(w).
  \]

  \item Use the fractional linear transform to take $\phi$ back to $u$ on the
  right half-plane.
\end{enumerate}

\begin{solution}
% Write your solution here.
\end{solution}

\section{Optional Warm-ups}

\subsection*{18.04 PS4, Problem 6}

\begin{enumerate}[(a)]
  \item Suppose that $f(z)$ is analytic on a region $A$ that contains the disk
  $|z-z_0|\le r$. Use Cauchy's integral formula to prove the mean value
  property
  \[
    f(z_0)=\frac{1}{2\pi}\int_0^{2\pi}
    f(z_0+re^{i\theta})\,\dd\theta.
  \]

  \item Prove the more general formula
  \[
    f^{(n)}(z_0)
    =
    \frac{n!}{2\pi r^n}
    \int_0^{2\pi}f(z_0+re^{i\theta})e^{-in\theta}\,\dd\theta.
  \]
\end{enumerate}

\begin{solution}
% Write your solution here.
\end{solution}

\subsection*{18.04 PS4, Problem 10}

Suppose $f(z)$ is entire and
\[
  \lim_{z\to\infty}\frac{f(z)}{z}=0.
\]
Show that $f(z)$ is constant.

You may use Morera's theorem: if $g(z)$ is analytic on $A-\{z_0\}$ and
continuous on $A$, then $g$ is analytic on $A$.

\begin{solution}
% Write your solution here.
\end{solution}

\section{Optional 18.112 References}

The following are useful harder references from MIT OCW
\href{https://ocw.mit.edu/courses/18-112-functions-of-a-complex-variable-fall-2008/pages/assignments/}{18.112 Functions of a Complex Variable, Fall 2008}.
They are assigned from Ahlfors, \emph{Complex Analysis}, 3rd edition, so this
file records the references rather than retyping the textbook problem
statements.

\begin{enumerate}
  \item Problem set 4: Ahlfors p.~130, Problem 3. Use it for essential
  singularities at infinity of $e^z$, $\sin z$, and $\cos z$.

  \item Problem set 5: Ahlfors p.~161, Problems 3(b), 3(f), and 3(h). Use them
  for residue-theorem definite integrals and logarithm branch choices.
\end{enumerate}

\end{document}
